% @see : https://coopmaths.fr/alea/?uuid=051aa&id=auto6G8A&n=1&d=10&s=1&s2=true&cd=1&alea=zndk&uuid=051aa&id=auto6G8A&n=1&d=10&s=2&s2=true&cd=1&alea=D0mb&uuid=051aa&id=auto6G8A&n=1&d=10&s=3&s2=true&cd=1&alea=xlS7&uuid=051aa&id=auto6G8A&n=1&d=10&s=5&s2=true&cd=1&alea=VfHD&uuid=051aa&id=auto6G8A&n=1&d=10&s=6&s2=true&cd=1&alea=hIH7&uuid=051aa&id=auto6G8A&n=1&d=10&s=7&s2=true&cd=1&alea=fOkp&uuid=5f117&id=6M3B&n=1&d=10&s2=1&s3=false&cd=1&alea=sXKn&uuid=e332d&id=6M3C-flash1&alea=X9Tn&pdfParam=eyJ0aXRsZSI6IiIsInJlZmVyZW5jZSI6IiIsInN1YnRpdGxlIjoiIiwic3R5bGUiOiJDb29wbWF0aHMiLCJmb250T3B0aW9uIjoiU3RhbmRhcmRGb250IiwidGFpbGxlRm9udE9wdGlvbiI6MTIsImR5c1RhaWxsZUZvbnRPcHRpb24iOjE0LCJjb3JyZWN0aW9uT3B0aW9uIjoiQXZlY0NvcnJlY3Rpb24iLCJxcmNvZGVPcHRpb24iOiJTYW5zUXJjb2RlIiwidHlwZUZpY2hlIjoiRmljaGUiLCJkdXJhdGlvbkNhbk9wdGlvbiI6IjkgbWluIiwidGl0bGVPcHRpb24iOiJTYW5zVGl0cmUiLCJuYlZlcnNpb25zIjoxLCJleG9zIjp7fX0\%3D&v=eleve
\documentclass[a4paper,11pt,fleqn]{article}

%%% COULEURS %%%
\usepackage[table,svgnames]{xcolor}
\definecolor{nombres}{cmyk}{0,.8,.95,0}
\definecolor{gestion}{cmyk}{.75,1,.11,.12}
\definecolor{gestionbis}{cmyk}{.75,1,.11,.12}
\definecolor{grandeurs}{cmyk}{.02,.44,1,0}
\definecolor{geo}{cmyk}{.62,.1,0,0}
\definecolor{algo}{cmyk}{.69,.02,.36,0}
\definecolor{correction}{cmyk}{.63,.23,.93,.06}
\arrayrulecolor{couleur_theme} % Couleur des filets des tableaux

\usepackage[most]{tcolorbox} % Supprimé à cause de conflits avec ProfCollege
\tcbuselibrary{documentation}

%%%%%%%% Modification paramétrage pour footnotes
% Le paquet semble déjà appelé en amont, curieux qu'il n'y ait pas eu de clash
% avec l'appel ci-dessous mais seulement quand j'ai essayé de l'inclure avec l'option lualatex

% \usepackage{hyperref}
% Parametrages
\hypersetup{
    colorlinks=true,% On active la couleur pour les liens. Couleur par défaut rouge
    linkcolor=black,% On définit la couleur pour les liens internes
    % filecolor=magenta,% On définit la couleur pour les liens vers les fichiers locaux      
    urlcolor=black,% On définit la couleur pour les liens vers des sites web
    % pdftitle={Puissance Quatre},% On définit un titre pour le document pdf
    % pdfpagemode=FullScreen,% On fixe l'affichage par défaut à plein écran
    }
% Pour avoir les numérotations arabic y compris dans les environnements tcolorbox
\renewcommand{\thempfootnote}{\arabic{mpfootnote}}
%%%%%%%%

\usepackage{multicol} 
\usepackage{calc}
\usepackage{enumitem}
\usepackage{graphicx}
\usepackage{tabularx}
\usepackage{tablvar}

\setlength{\parindent}{0mm}
\renewcommand{\arraystretch}{1.5}
\renewcommand{\labelenumi}{\textbf{\theenumi{}.}}
\renewcommand{\labelenumii}{\textbf{\theenumii{}.}}
\setlength{\fboxsep}{3mm}

\setlength{\headheight}{14.5pt}

\setlength{\premulticols}{0pt}

\newcounter{ExoMA}

\newlength{\parindentMA}%
\setlength{\parindentMA}{\parindent}
\addtolength{\parindentMA}{30pt}

\usepackage{simplekv}

\setKVdefault[Theme]{Coopmath,Classiques=false,CANs=false}
\defKV[Theme]{Classique=\setKV[Theme]{Coopmath=false}\setKV[Theme]{Classiques}}
\defKV[Theme]{CAN=\setKV[Theme]{Coopmath=false}\setKV[Theme]{CANs}}

\usepackage{fancyhdr}
\fancyhead[C]{}
\fancyfoot{}

\def\bla{}

\NewDocumentCommand\Theme{o m m m m}{%
  \useKVdefault[Theme]%
  \setKV[Theme]{#1}%
  \ifboolKV[Theme]{CANs}{%
    \usepackage[left=1.5cm,right=1.5cm,top=2cm,bottom=2cm]{geometry}%
    \fancypagestyle{premierePage}{%
      \fancyhead[C]{\textsc{\ifx\bla#2\bla Course aux nombres\else#2\fi}}%
      \fancyfoot[R]{\scriptsize Coopmaths.fr -- CC-BY-SA}%
      \setlength{\headheight}{14.5pt}%
      \NewDocumentCommand\dureeCan{}{\ifx\bla#4\bla 9 minutes\else #4\fi}
      \NewDocumentCommand\titreSujetCan{}{\ifx\bla#3\bla\else\textbf{#3}\fi}
    }%
    \pagestyle{premierePage}
    \colorlet{couleur_theme}{black}%
    \colorlet{couleur_numerotation}{couleur_theme}%
    \let\EXO\EXOCan\let\endEXO\endEXOCan%
  }{%
    \renewcommand{\headrulewidth}{0pt}
    \renewcommand{\footrulewidth}{0pt}
    \pagestyle{fancy}
    \ifboolKV[Theme]{Classiques}{%
      \usepackage[left=0.65cm,right=0.65cm,top=2cm,bottom=1.5cm]{geometry}
      \let\EXO\EXOlibre\let\endEXO\endEXOlibre%
      \colorlet{couleur_theme}{black}%
      \colorlet{couleur_numerotation}{couleur_theme}%
      \fancyhead[C]{\bfseries #3}
      \fancyhead[L]{\bfseries #4}
      \fancyfoot[R]{#5}
    }{%
      \usepackage[left=1.5cm,right=1.5cm,top=4cm,bottom=2cm]{geometry}
      \theme{#2}{#3}{#4}{#5}%
      \let\EXO\EXOcoop\let\endEXO\endEXOcoop%
    }%
  }%
}%

\NewDocumentEnvironment{EXOCan}{m m +b}{%
  #3
}

\NewDocumentEnvironment{EXOcoop}{m m +b}{%
  \begin{tcolorbox}[%
    enhanced,
    breakable,
    colback=white,
    colframe=white,
    coltitle=black,
    title=\IfNoValueTF{#1}{}{\textmd{#1}},%
    overlay unbroken and first={%
      \IfNoValueTF{#2}{}{%
        \node[%
        fill=white,
        anchor=east,
        xshift=10pt,
        text=black
        ]
        at (frame.north east)
        {\scriptsize #2};
      }
      \node[anchor=west,xshift=-20pt,yshift=-15pt] at (frame.north west) {\stepcounter{ExoMA}\numb{\arabic{ExoMA}}};
    }
    ]
    #3
  \end{tcolorbox}
}

\NewDocumentEnvironment{EXOlibre}{m m +b}{%
  \begin{tcolorbox}[%
    enhanced,
    breakable,
    colback=white,
    colframe=white,%couleur_theme,
    coltitle=black,
    title=\stepcounter{ExoMA}\textbf{Exercice \arabic{ExoMA}},%
    ]
    \IfNoValueTF{#1}{}{#1\par}%
    #3
  \end{tcolorbox}%
}


%%% MISE EN PAGE %%%
\usepackage{setspace}
\usepackage{pgf}%
\usepackage{pgfplots}
\pgfplotsset{compat=1.18}
\usetikzlibrary{babel,arrows, arrows.meta,calc,fit,patterns,plotmarks,shapes.geometric,shapes.misc,shapes.symbols,shapes.arrows,shapes.callouts, shapes.multipart, shapes.gates.logic.US,shapes.gates.logic.IEC, er, automata,backgrounds,chains,topaths,trees,petri,mindmap,matrix, calendar, folding,fadings,through,positioning,scopes,decorations.fractals,decorations.shapes,decorations.text,decorations.pathmorphing,decorations.pathreplacing,decorations.footprints,decorations.markings,shadows}

\spaceskip=2\fontdimen2\font plus 3\fontdimen3\font minus3\fontdimen4\font\relax %Pour doubler l'espace entre les mots
\newcommand\numb[1]{ % Dessin autour du numéro d'exercice
  \begin{tikzpicture}[overlay,scale=.8]%,yshift=-.3cm
    \draw[fill=couleur_numerotation,couleur_numerotation](-.4,-0.1)rectangle($(-0.4,.8)+(2em,0)$);
    \draw (-.4,.8) node[white,anchor=north west,inner sep=0.5pt]{\bfseries EX}; 
    \draw[line width=.05cm,couleur_numerotation,fill=white] (0,0)--(.5,.5)--(1,0)--(.5,-.5)--cycle;
    \draw (0.5,0) node[anchor=center,couleur_numerotation]{\large\bfseries#1};
  \end{tikzpicture}
}

% echelle pour le dé
\def\globalscale {0.04}
% abscisse initiale pour les chevrons
\def\xini {3}

\newcommand\theme[4]{%
  \fancyhead[C]{%
    % Tracé du dé
    \begin{tikzpicture}[y=0.80pt, x=0.80pt, yscale=-\globalscale, xscale=\globalscale,remember picture, overlay,transform canvas={xshift=-0.5\linewidth,yshift=2.7cm},fill=couleur_theme]%,xshift=17cm,yshift=9.5cm
      %%%% Arc supérieur gauche%%%%
      \path[fill](523,1424)..controls(474,1413)and(404,1372)..(362,1333)..controls(322,1295)and(313,1272)..(331,1254)..controls(348,1236)and(369,1245)..(410,1283)..controls(458,1328)and(517,1356)..(575,1362)..controls(635,1368)and(646,1375)..(643,1404)..controls(641,1428)and(641,1428)..(596,1430)..controls(571,1431)and(538,1428)..(523,1424)--cycle;
      %%%% Dé face supérieur%%%%
      \path[fill](512,1272)..controls(490,1260)and(195,878)..(195,861)..controls(195,854)and(198,846)..(202,843)..controls(210,838)and(677,772)..(707,772)..controls(720,772)and(737,781)..(753,796)..controls(792,833)and(1057,1179)..(1057,1193)..controls(1057,1200)and(1053,1209)..(1048,1212)..controls(1038,1220)and(590,1283)..(551,1282)..controls(539,1282)and(521,1278)..(512,1272)--cycle;
      %%%% Dé faces gauche et droite%%%%
      \path[fill](1061,1167)..controls(1050,1158)and(978,1068)..(900,967)..controls(792,829)and(756,777)..(753,756)--(748,729)--(724,745)..controls(704,759)and(660,767)..(456,794)..controls(322,813)and(207,825)..(200,822)..controls(193,820)and(187,812)..(187,804)..controls(188,797)and(229,688)..(279,563)..controls(349,390)and(376,331)..(391,320)..controls(406,309)and(462,299)..(649,273)..controls(780,254)and(897,240)..(907,241)..controls(918,243)and(927,249)..(928,256)..controls(930,264)and(912,315)..(889,372)..controls(866,429)and(848,476)..(849,477)..controls(851,479)and(872,432)..(897,373)..controls(936,276)and(942,266)..(960,266)..controls(975,266)and(999,292)..(1089,408)..controls(1281,654)and(1290,666)..(1290,691)..controls(1290,720)and(1104,1175)..(1090,1180)..controls(1085,1182)and (1071,1176)..(1061,1167)--cycle;
      %%%% Arc inférieur bas%%%%
      \path[fill](1329,861)..controls(1316,848)and(1317,844)..(1339,788)..controls(1364,726)and(1367,654)..(1347,591)..controls(1330,539)and(1338,522)..(1375,526)..controls(1395,528)and(1400,533)..(1412,566)..controls(1432,624)and(1426,760)..(1401,821)..controls(1386,861)and(1380,866)..(1361,868)..controls(1348,870)and(1334,866)..(1329,861)--cycle;
      %%%% Arc inférieur gauche%%%%
      \path[fill](196,373)..controls(181,358)and(186,335)..(213,294)..controls(252,237)and(304,190)..(363,161)..controls(435,124)and(472,127)..(472,170)..controls(472,183)and(462,192)..(414,213)..controls(350,243)and(303,283)..(264,343)..controls(239,383)and(216,393)..(196,373)--cycle;
    \end{tikzpicture}
    \begin{tikzpicture}[line cap=round,line join=round,remember picture, overlay, shift={(current page.north west)},yshift=-8.5cm]
      \fill[fill=couleur_theme] (0,5) rectangle (21,6);
      \fill[fill=couleur_theme] (\xini,6)--(\xini+1.5,6)--(\xini+2.5,7)--(\xini+1.5,8)--(\xini,8)--(\xini+1,7)-- cycle;
      \fill[fill=couleur_theme] (\xini+2,6)--(\xini+2.5,6)--(\xini+3.5,7)--(\xini+2.5,8)--(\xini+2,8)--(\xini+3,7)-- cycle;  
      \fill[fill=couleur_theme] (\xini+3,6)--(\xini+3.5,6)--(\xini+4.5,7)--(\xini+3.5,8)--(\xini+3,8)--(\xini+4,7)-- cycle;   
      \node[color=white] at (10.5,5.5) {\LARGE \bfseries{ \MakeUppercase{#4}}};
    \end{tikzpicture}
    \begin{tikzpicture}[remember picture,overlay]
      \node[anchor=north east,inner sep=0pt] at ($(current page.north east)+(0,-.8cm)$) {};
      \node[anchor=west, fill=white, inner sep = 0pt] at ($(current page.north west)+(0.2,-2.3cm)$) {\footnotesize \bfseries{MathALÉA}};
    \end{tikzpicture}
    \begin{tikzpicture}[remember picture,overlay]
      \node[anchor=north east,inner sep=0pt] at ($(current page.north east)+(0,-.8cm)$) {};
      \node[anchor=east, fill=white] at ($(current page.north east)+(-2,-1.5cm)$) {\Huge \textcolor{couleur_theme}{\bfseries{\#}} \bfseries{#2} \textcolor{couleur_theme}{\bfseries \MakeUppercase{#3}}};
    \end{tikzpicture}
  }%
  \fancyfoot[R]{%
    \begin{tikzpicture}[remember picture,overlay]
      \node[anchor=south east] at ($(current page.south east)+(-2,0.25cm)$) {\scriptsize {\bfseries \href{https://coopmaths.fr/}{Coopmaths.fr} -- \href{http://creativecommons.fr/licences/}{CC-BY-SA}}};
    \end{tikzpicture}
    \begin{tikzpicture}[line cap=round,line join=round,remember picture, overlay,xscale=0.5,yscale=0.5, shift={(current page.south west)},xshift=35.7cm,yshift=-6cm]
      \fill[fill=couleur_theme] (\xini,6)--(\xini+1.5,6)--(\xini+2.5,7)--(\xini+1.5,8)--(\xini,8)--(\xini+1,7)-- cycle;
      \fill[fill=couleur_theme] (\xini+2,6)--(\xini+2.5,6)--(\xini+3.5,7)--(\xini+2.5,8)--(\xini+2,8)--(\xini+3,7)-- cycle;  
      \fill[fill=couleur_theme] (\xini+3,6)--(\xini+3.5,6)--(\xini+4.5,7)--(\xini+3.5,8)--(\xini+3,8)--(\xini+4,7)-- cycle;  
    \end{tikzpicture}
  }
  \fancyfoot[C]{}
  \colorlet{couleur_theme}{#1}
  \colorlet{couleur_numerotation}{couleur_theme}
  \def\iconeobjectif{icone-objectif-#1}
  \def\urliconeomethode{icone-methode-#1}
}%

\newcommand*\tikzfootMA{%
  \ifboolKV[Theme]{CANs}{
    \begin{tikzpicture}[remember picture,overlay,shift=(current page.north west)]
      \begin{scope}[x=\textwidth-\oddsidemargin,y=\textheight+5pt]
        \draw[correction,line width=4pt,dashed,dash pattern= on 10pt off 10pt,shift={(-\oddsidemargin,\topmargin)}] (0,0) rectangle (1,-1);
      \end{scope}
    \end{tikzpicture} 
  }{%
  \ifboolKV[Theme]{Classiques}{%
  \begin{tikzpicture}[remember picture,overlay,shift=(current page.south west)]
    \begin{scope}[x={(current page.south east)},y={(current page.north west)}]
      \draw[correction,line width=4pt,dashed,dash pattern= on 10pt off 10pt] ($(0,0)+(5mm,1cm)$) rectangle ($(1,1)+(-5mm,-1.75cm)$);
    \end{scope}
  \end{tikzpicture} 
  }{%
  \begin{tikzpicture}[remember picture,overlay,shift=(current page.south west)]
    \begin{scope}[x={(current page.south east)},y={(current page.north west)}]
      \draw[correction,line width=4pt,dashed,dash pattern= on 10pt off 10pt] ($(0,0)+(5mm,1.5cm)$) rectangle ($(1,1)+(-5mm,-3.75cm)$);
    \end{scope}
  \end{tikzpicture} 
  }%
  }
}

\newenvironment{Correction}{%
  \setcounter{ExoMA}{0}%
  \cfoot{\tikzfootMA}
}{}%

\newcommand\version[1]{
  \fancyhead[R]{
    \begin{tikzpicture}[remember picture,overlay]
    \node[anchor=north east,inner sep=0pt] at ($(current page.north east)+(-.5,-.5cm)$) {\large \textcolor{couleur_theme}{\bfseries V#1}};
    \end{tikzpicture}
  }
}

\usepackage[french]{babel}
% loadPackagesFromContent
\usepackage{tikz}
\newcommand{\e}{\mathrm{e}}
\usepackage{amssymb}
\usepackage{etoolbox}
\newbool{dys}
\setbool{dys}{false}          
\ifbool{dys}{
% POLICE DYS
% ===== VARIABLE =====
\def\FontChoisie{Fira} % valeurs possibles : Fira, lmodern, tgheros
\newcommand{\choiceFontsDys}[1]{
\ifstrequal{#1}{Fira}{%
  % Fira Sans + Fira Math
  % Description : Fira Sans est une police moderne, et Fira Math est une version mathématique compatible qui maintient le style sans-serif.
  % Utilisation : Fonctionne bien avec XeLaTeX et LuaLaTeX.
  \usepackage{fontspec}
  \setmainfont{Fira Sans}
  \setsansfont{Fira Sans}
  \usepackage{unicode-math}
  \setmathfont{Fira Math}
}{}
\ifstrequal{#1}{lmodern}{%
  % Latin Modern Sans
  % Description : Une version sans-serif de la célèbre police Latin Modern, qui est compatible avec mathastext.
  % Utilisation : Fonctionne avec pdfLaTeX, XeLaTeX et LuaLaTeX.
  \usepackage{lmodern}
  \renewcommand{\familydefault}{\sfdefault}
  \usepackage{mathastext}
}{}
\ifstrequal{#1}{tgheros}{%
  % TeX Gyre Heros + mathastext
  % Description : TeX Gyre Heros est une alternative à Helvetica, et le package mathastext permet d'utiliser la même police pour les mathématiques.
  % Utilisation : Fonctionne avec pdfLaTeX, XeLaTeX, ou LuaLaTeX.
  \usepackage{tgheros}
  \renewcommand{\familydefault}{\sfdefault}
  \usepackage{mathastext}
}{}
}
% Fira ou lmodern ou tgheros
\choiceFontsDys{\FontChoisie}

%\usepackage{unicode-math}
%\usepackage{fontspec}
%\setmainfont{TeX Gyre Schola}
%\setsansfont{TeX Gyre Schola}
%\setmathfont{TeX Gyre Schola Math}
%\setmainfont{OpenDyslexic}[Scale=1.0]
%\setsansfont{OpenDyslexic}[Scale=1.0]
%\setmathfont{OpenDyslexic}[Scale=1.0,range=up/{Latin,latin,num}]
\usepackage[fontsize=14]{scrextend}
\usepackage{setspace}
\setstretch{1.7}
% Une valeur d'environ 1.2 à 1.7 est souvent recommandée. Cela permet d'augmenter l'espace entre les lignes tout en offrant un léger espacement entre les mots.
\setlength{\spaceskip}{1.2em}  
% Une valeur d'environ 1.2em à 1.5em est couramment conseillée. Cela crée un espace plus ample entre les mots, ce qui peut aider à réduire la fatigue visuelle et à améliorer la fluidité de la lecture.
}{
% POLICE STANDARD
\usepackage[T1]{fontenc} 
\usepackage[scaled=1]{helvet}
\usepackage[fontsize=12]{scrextend}
}

\Theme[Coopmaths]{nombres}{}{}{}
\begin{document}

% @see : https://coopmaths.fr/alea?uuid=051aa&id=auto6G8A&n=1&d=10&s=1&s2=true&alea=zndk&cd=1&cols=4
\begin{EXO}{Donner le nom de ce solide.}{auto6G8A}
\begin{multicols}{4}
 \begin{minipage}[t]{\linewidth} \begin{tikzpicture}[baseline,scale = 0.5]

    \tikzset{
      point/.style={
        thick,
        draw,
        cross out,
        inner sep=0pt,
        minimum width=5pt,
        minimum height=5pt,
      },
    }
    \clip (-2.5,-1) rectangle (2.67,3.93);
    	\draw[color ={black}] (2.17,0.25)--(1.75,0.43);
	\draw[color ={black}] (1.75,0.43)--(-2,0);
	\draw[color ={black}] (-2,0)--(-0.87,-0.5);
	\draw[color ={black}] (-0.87,-0.5)--(2.17,0.25);
	\draw[color ={black}] (2.17,3.25)--(1.75,3.43);
	\draw[color ={black}] (1.75,3.43)--(-2,3);
	\draw[color ={black}, densely dash dot dot ] (-2,3)--(-0.87,2.5);
	\draw[color ={black}, densely dash dot dot ] (-0.87,2.5)--(2.17,3.25);
	\draw[color ={black}] (2.17,0.25)--(2.17,3.25);
	\draw[color ={black}] (1.75,0.43)--(1.75,3.43);
	\draw[color ={black}] (-2,0)--(-2,3);
	\draw[color ={black}, densely dash dot dot ] (-0.87,-0.5)--(-0.87,2.5);

\end{tikzpicture} 
    
    	$\square\;$ Cône\qquad $\square\;$ Pyramide\qquad $\square\;$ Cylindre\qquad $\square\;$ Prisme\qquad $\square\;$ Sphère\qquad $\square\;$ Pavé\qquad $\square\;$ Cube\qquad \\
 \end{minipage}
\end{multicols}
\end{EXO}

% @see : https://coopmaths.fr/alea?uuid=051aa&id=auto6G8A&n=1&d=10&s=2&s2=true&alea=D0mb&cd=1&cols=4
\begin{EXO}{Donner le nom de ce solide.}{auto6G8A}
\begin{multicols}{4}
 \begin{minipage}[t]{\linewidth} \begin{tikzpicture}[baseline,scale = 0.5]

    \tikzset{
      point/.style={
        thick,
        draw,
        cross out,
        inner sep=0pt,
        minimum width=5pt,
        minimum height=5pt,
      },
    }
    \clip (-2.63,-0.98) rectangle (2.5,3.5);
    	\draw[color ={black}] (2,0)--(0,3);
	\draw[color ={black}, densely dash dot dot ] (1.44,0.48)--(0,3);
	\draw[color ={black}, densely dash dot dot ] (-1.11,0.29)--(0,3);
	\draw[color ={black}] (-2.13,-0.29)--(0,3);
	\draw[color ={black}] (-0.21,-0.48)--(0,3);
	\draw[color ={black}, densely dash dot dot ] (2,0)--(1.44,0.48);
	\draw[color ={black}, densely dash dot dot ] (1.44,0.48)--(-1.11,0.29);
	\draw[color ={black}, densely dash dot dot ] (-1.11,0.29)--(-2.13,-0.29);
	\draw[color ={black}] (-2.13,-0.29)--(-0.21,-0.48);
	\draw[color ={black}] (-0.21,-0.48)--(2,0);

\end{tikzpicture} 
    
    	$\square\;$ Sphère\qquad $\square\;$ Pavé\qquad $\square\;$ Cube\qquad $\square\;$ Cylindre\qquad $\square\;$ Pyramide\qquad $\square\;$ Prisme\qquad $\square\;$ Cône\qquad \\
 \end{minipage}
\end{multicols}
\end{EXO}

% @see : https://coopmaths.fr/alea?uuid=051aa&id=auto6G8A&n=1&d=10&s=3&s2=true&alea=xlS7&cd=1&cols=4
\begin{EXO}{Donner le nom de ce solide.}{auto6G8A}
\begin{multicols}{4}
 \begin{minipage}[t]{\linewidth} \begin{tikzpicture}[baseline,scale = 0.5]

    \tikzset{
      point/.style={
        thick,
        draw,
        cross out,
        inner sep=0pt,
        minimum width=5pt,
        minimum height=5pt,
      },
    }
    \clip (-2.9,-1.3666666666666667) rectangle (2.9,4.25);
    	\draw[color={black}, dash dot ] (2.3,0) arc [start angle=0, end angle = 180, x radius = 2.3, y radius = 0.7666666666666666];
	\draw[color={black}] (2.3,0) arc [start angle=0, end angle = -180, x radius = 2.3, y radius = 0.7666666666666666];
	\draw[color ={black}] (2.3,0.1)--(0,3);
	\draw[color ={black}] (-2.3,0.1)--(0,3);
	\draw[color ={black}, densely dash dot dot ] (0,-0.75)--(0,3.75);
	\draw[color ={black},line width = 0.625,opacity = 0.8] (-0.19,0.19)--(0.19,-0.19);\draw[color ={black},line width = 0.625,opacity = 0.8] (-0.19,-0.19)--(0.19,0.19);

\end{tikzpicture} 
    
    	$\square\;$ Sphère\qquad $\square\;$ Pavé\qquad $\square\;$ Cylindre\qquad $\square\;$ Cône\qquad $\square\;$ Prisme\qquad $\square\;$ Pyramide\qquad $\square\;$ Cube\qquad \\
 \end{minipage}
\end{multicols}
\end{EXO}

% @see : https://coopmaths.fr/alea?uuid=051aa&id=auto6G8A&n=1&d=10&s=5&s2=true&alea=VfHD&cd=1&cols=4
\begin{EXO}{Donner le nom de ce solide.}{auto6G8A}
\begin{multicols}{4}
 \begin{minipage}[t]{\linewidth} \begin{tikzpicture}[baseline,scale = 0.5]

    \tikzset{
      point/.style={
        thick,
        draw,
        cross out,
        inner sep=0pt,
        minimum width=5pt,
        minimum height=5pt,
      },
    }
    \clip (-2.1100000000000003,-2.27) rectangle (5.11,2.27);
    	\draw[color ={black}] (-0.39,-1.06)--(-0.39,1.77);
	\draw[color ={black}] (-0.39,1.77)--(-1.61,1.06);
	\draw[color ={black}] (-1.61,1.06)--(-1.61,-1.77);
	\draw[color ={black}] (-1.61,-1.77)--(-0.39,-1.06);
	\draw[color ={black}] (4.61,-1.06)--(4.61,1.77);
	\draw[color ={black}, densely dash dot dot ] (4.61,1.77)--(3.39,1.06);
	\draw[color ={black}, densely dash dot dot ] (3.39,1.06)--(3.39,-1.77);
	\draw[color ={black}] (3.39,-1.77)--(4.61,-1.06);
	\draw[color ={black}] (-0.39,-1.06)--(4.61,-1.06);
	\draw[color ={black}] (-0.39,1.77)--(4.61,1.77);
	\draw[color ={black}, densely dash dot dot ] (-1.61,1.06)--(3.39,1.06);
	\draw[color ={black}] (-1.61,-1.77)--(3.39,-1.77);

\end{tikzpicture} 
    
    	$\square\;$ Cube\qquad $\square\;$ Cône\qquad $\square\;$ Cylindre\qquad $\square\;$ Sphère\qquad $\square\;$ Prisme\qquad $\square\;$ Pavé\qquad $\square\;$ Pyramide\qquad \\
 \end{minipage}
\end{multicols}
\end{EXO}

% @see : https://coopmaths.fr/alea?uuid=051aa&id=auto6G8A&n=1&d=10&s=6&s2=true&alea=hIH7&cd=1&cols=4
\begin{EXO}{Donner le nom de ce solide.}{auto6G8A}
\begin{multicols}{4}
 \begin{minipage}[t]{\linewidth} \begin{tikzpicture}[baseline,scale = 0.5]

    \tikzset{
      point/.style={
        thick,
        draw,
        cross out,
        inner sep=0pt,
        minimum width=5pt,
        minimum height=5pt,
      },
    }
    \clip (-2.1100000000000003,-2.27) rectangle (3.11,2.27);
    	\draw[color ={black}] (-0.39,-1.06)--(-0.39,1.77);
	\draw[color ={black}] (-0.39,1.77)--(-1.61,1.06);
	\draw[color ={black}] (-1.61,1.06)--(-1.61,-1.77);
	\draw[color ={black}] (-1.61,-1.77)--(-0.39,-1.06);
	\draw[color ={black}] (2.61,-1.06)--(2.61,1.77);
	\draw[color ={black}, densely dash dot dot ] (2.61,1.77)--(1.39,1.06);
	\draw[color ={black}, densely dash dot dot ] (1.39,1.06)--(1.39,-1.77);
	\draw[color ={black}] (1.39,-1.77)--(2.61,-1.06);
	\draw[color ={black}] (-0.39,-1.06)--(2.61,-1.06);
	\draw[color ={black}] (-0.39,1.77)--(2.61,1.77);
	\draw[color ={black}, densely dash dot dot ] (-1.61,1.06)--(1.39,1.06);
	\draw[color ={black}] (-1.61,-1.77)--(1.39,-1.77);

\end{tikzpicture} 
    
    	$\square\;$ Sphère\qquad $\square\;$ Pavé\qquad $\square\;$ Cône\qquad $\square\;$ Pyramide\qquad $\square\;$ Cylindre\qquad $\square\;$ Cube\qquad $\square\;$ Prisme\qquad \\
 \end{minipage}
\end{multicols}
\end{EXO}

% @see : https://coopmaths.fr/alea?uuid=051aa&id=auto6G8A&n=1&d=10&s=7&s2=true&alea=fOkp&cd=1&cols=4
\begin{EXO}{Donner le nom de ce solide.}{auto6G8A}
\begin{multicols}{4}
 \begin{minipage}[t]{\linewidth} \begin{tikzpicture}[baseline,scale = 0.5]

    \tikzset{
      point/.style={
        thick,
        draw,
        cross out,
        inner sep=0pt,
        minimum width=5pt,
        minimum height=5pt,
      },
    }
    \clip (-2.6,-2.5) rectangle (2.6,2.5);
    	\draw[color={black}] (0,0) circle (2);
	\draw[color={black}, dash dot ] (2,0) arc [start angle=0, end angle = 180, x radius = 2, y radius = 0.6666666666666666];
	\draw[color={black}] (2,0) arc [start angle=0, end angle = -180, x radius = 2, y radius = 0.6666666666666666];

\end{tikzpicture} 
    
    	$\square\;$ Sphère\qquad $\square\;$ Cube\qquad $\square\;$ Prisme\qquad $\square\;$ Cylindre\qquad $\square\;$ Pyramide\qquad $\square\;$ Pavé\qquad $\square\;$ Cône\qquad \\
 \end{minipage}
\end{multicols}
\end{EXO}

% @see : https://coopmaths.fr/alea?uuid=5f117&id=6M3B&n=1&d=10&s2=1&s3=false&alea=sXKn&cd=1&cols=1
\begin{EXO}{}{6M3B}

 Ci-dessous, un empilement de cubes est représenté sous deux angles différents et un pavé droit est représenté à côté.\\\begin{tikzpicture}[baseline,scale = 0.6]

    \tikzset{
      point/.style={
        thick,
        draw,
        cross out,
        inner sep=0pt,
        minimum width=5pt,
        minimum height=5pt,
      },
    }
    \clip (-1.0209445330007911,-0.5) rectangle (3.2807750813696934,5.0320566851822495);
    	\draw[color={black},preaction={fill,color = {green}}] (1.622319,1.489282)--(2.607127,1.600901)--(2.607127,2.366946)--(1.622319,2.255327)--cycle;
	\draw[color={black},preaction={fill,color = {white}}] (2.607127,2.366946)--(2.433479,2.999968)--(1.448671,2.888349)--(1.622319,2.255327)--cycle;
	\draw[color={black},preaction={fill,color = {gray}}] (1.622319,2.255327)--(1.448671,2.888349)--(1.448671,2.122304)--(1.622319,1.489282)--cycle;
	\draw[color={black},preaction={fill,color = {green}}] (1.622319,2.255327)--(2.607127,2.366946)--(2.607127,3.13299)--(1.622319,3.021371)--cycle;
	\draw[color={black},preaction={fill,color = {white}}] (2.607127,3.13299)--(2.433479,3.766012)--(1.448671,3.654393)--(1.622319,3.021371)--cycle;
	\draw[color={black},preaction={fill,color = {gray}}] (1.622319,3.021371)--(1.448671,3.654393)--(1.448671,2.888349)--(1.622319,2.255327)--cycle;
	\draw[color={black},preaction={fill,color = {green}}] (1.622319,3.021371)--(2.607127,3.13299)--(2.607127,3.899034)--(1.622319,3.787416)--cycle;
	\draw[color={black},preaction={fill,color = {white}}] (2.607127,3.899034)--(2.433479,4.532057)--(1.448671,4.420438)--(1.622319,3.787416)--cycle;
	\draw[color={black},preaction={fill,color = {gray}}] (1.622319,3.787416)--(1.448671,4.420438)--(1.448671,3.654393)--(1.622319,3.021371)--cycle;
	\draw[color={black},preaction={fill,color = {green}}] (1.795967,0.85626)--(2.780775,0.967879)--(2.780775,1.733923)--(1.795967,1.622304)--cycle;
	\draw[color={black},preaction={fill,color = {white}}] (2.780775,1.733923)--(2.607127,2.366946)--(1.622319,2.255327)--(1.795967,1.622304)--cycle;
	\draw[color={black},preaction={fill,color = {gray}}] (1.795967,1.622304)--(1.622319,2.255327)--(1.622319,1.489282)--(1.795967,0.85626)--cycle;
	\draw[color={black},preaction={fill,color = {green}}] (1.795967,1.622304)--(2.780775,1.733923)--(2.780775,2.499968)--(1.795967,2.388349)--cycle;
	\draw[color={black},preaction={fill,color = {white}}] (2.780775,2.499968)--(2.607127,3.13299)--(1.622319,3.021371)--(1.795967,2.388349)--cycle;
	\draw[color={black},preaction={fill,color = {gray}}] (1.795967,2.388349)--(1.622319,3.021371)--(1.622319,2.255327)--(1.795967,1.622304)--cycle;
	\draw[color={black},preaction={fill,color = {green}}] (0.637511,1.377663)--(1.622319,1.489282)--(1.622319,2.255327)--(0.637511,2.143708)--cycle;
	\draw[color={black},preaction={fill,color = {white}}] (1.622319,2.255327)--(1.448671,2.888349)--(0.463863,2.77673)--(0.637511,2.143708)--cycle;
	\draw[color={black},preaction={fill,color = {gray}}] (0.637511,2.143708)--(0.463863,2.77673)--(0.463863,2.010686)--(0.637511,1.377663)--cycle;
	\draw[color={black},preaction={fill,color = {green}}] (0.637511,2.143708)--(1.622319,2.255327)--(1.622319,3.021371)--(0.637511,2.909752)--cycle;
	\draw[color={black},preaction={fill,color = {white}}] (1.622319,3.021371)--(1.448671,3.654393)--(0.463863,3.542774)--(0.637511,2.909752)--cycle;
	\draw[color={black},preaction={fill,color = {gray}}] (0.637511,2.909752)--(0.463863,3.542774)--(0.463863,2.77673)--(0.637511,2.143708)--cycle;
	\draw[color={black},preaction={fill,color = {green}}] (0.637511,2.909752)--(1.622319,3.021371)--(1.622319,3.787416)--(0.637511,3.675797)--cycle;
	\draw[color={black},preaction={fill,color = {white}}] (1.622319,3.787416)--(1.448671,4.420438)--(0.463863,4.308819)--(0.637511,3.675797)--cycle;
	\draw[color={black},preaction={fill,color = {gray}}] (0.637511,3.675797)--(0.463863,4.308819)--(0.463863,3.542774)--(0.637511,2.909752)--cycle;
	\draw[color={black},preaction={fill,color = {green}}] (0.81116,0.744641)--(1.795967,0.85626)--(1.795967,1.622304)--(0.81116,1.510686)--cycle;
	\draw[color={black},preaction={fill,color = {white}}] (1.795967,1.622304)--(1.622319,2.255327)--(0.637511,2.143708)--(0.81116,1.510686)--cycle;
	\draw[color={black},preaction={fill,color = {gray}}] (0.81116,1.510686)--(0.637511,2.143708)--(0.637511,1.377663)--(0.81116,0.744641)--cycle;
	\draw[color={black},preaction={fill,color = {green}}] (0.81116,1.510686)--(1.795967,1.622304)--(1.795967,2.388349)--(0.81116,2.27673)--cycle;
	\draw[color={black},preaction={fill,color = {white}}] (1.795967,2.388349)--(1.622319,3.021371)--(0.637511,2.909752)--(0.81116,2.27673)--cycle;
	\draw[color={black},preaction={fill,color = {gray}}] (0.81116,2.27673)--(0.637511,2.909752)--(0.637511,2.143708)--(0.81116,1.510686)--cycle;
	\draw[color={black},preaction={fill,color = {green}}] (-0.347296,1.266044)--(0.637511,1.377663)--(0.637511,2.143708)--(-0.347296,2.032089)--cycle;
	\draw[color={black},preaction={fill,color = {white}}] (0.637511,2.143708)--(0.463863,2.77673)--(-0.520945,2.665111)--(-0.347296,2.032089)--cycle;
	\draw[color={black},preaction={fill,color = {gray}}] (-0.347296,2.032089)--(-0.520945,2.665111)--(-0.520945,1.899067)--(-0.347296,1.266044)--cycle;
	\draw[color={black},preaction={fill,color = {green}}] (-0.347296,2.032089)--(0.637511,2.143708)--(0.637511,2.909752)--(-0.347296,2.798133)--cycle;
	\draw[color={black},preaction={fill,color = {white}}] (0.637511,2.909752)--(0.463863,3.542774)--(-0.520945,3.431156)--(-0.347296,2.798133)--cycle;
	\draw[color={black},preaction={fill,color = {gray}}] (-0.347296,2.798133)--(-0.520945,3.431156)--(-0.520945,2.665111)--(-0.347296,2.032089)--cycle;
	\draw[color={black},preaction={fill,color = {green}}] (-0.347296,2.798133)--(0.637511,2.909752)--(0.637511,3.675797)--(-0.347296,3.564178)--cycle;
	\draw[color={black},preaction={fill,color = {white}}] (0.637511,3.675797)--(0.463863,4.308819)--(-0.520945,4.1972)--(-0.347296,3.564178)--cycle;
	\draw[color={black},preaction={fill,color = {gray}}] (-0.347296,3.564178)--(-0.520945,4.1972)--(-0.520945,3.431156)--(-0.347296,2.798133)--cycle;
	\draw[color={black},preaction={fill,color = {green}}] (-0.173648,0.633022)--(0.81116,0.744641)--(0.81116,1.510686)--(-0.173648,1.399067)--cycle;
	\draw[color={black},preaction={fill,color = {white}}] (0.81116,1.510686)--(0.637511,2.143708)--(-0.347296,2.032089)--(-0.173648,1.399067)--cycle;
	\draw[color={black},preaction={fill,color = {gray}}] (-0.173648,1.399067)--(-0.347296,2.032089)--(-0.347296,1.266044)--(-0.173648,0.633022)--cycle;
	\draw[color={black},preaction={fill,color = {green}}] (-0.173648,1.399067)--(0.81116,1.510686)--(0.81116,2.27673)--(-0.173648,2.165111)--cycle;
	\draw[color={black},preaction={fill,color = {white}}] (0.81116,2.27673)--(0.637511,2.909752)--(-0.347296,2.798133)--(-0.173648,2.165111)--cycle;
	\draw[color={black},preaction={fill,color = {gray}}] (-0.173648,2.165111)--(-0.347296,2.798133)--(-0.347296,2.032089)--(-0.173648,1.399067)--cycle;
	\draw[color={black},preaction={fill,color = {green}}] (0,0)--(0.984808,0.111619)--(0.984808,0.877663)--(0,0.766044)--cycle;
	\draw[color={black},preaction={fill,color = {white}}] (0.984808,0.877663)--(0.81116,1.510686)--(-0.173648,1.399067)--(0,0.766044)--cycle;
	\draw[color={black},preaction={fill,color = {gray}}] (0,0.766044)--(-0.173648,1.399067)--(-0.173648,0.633022)--(0,0)--cycle;

\end{tikzpicture}\begin{tikzpicture}[baseline,scale = 0.6]

    \tikzset{
      point/.style={
        thick,
        draw,
        cross out,
        inner sep=0pt,
        minimum width=5pt,
        minimum height=5pt,
      },
    }
    \clip (-2.220729309053138,-0.5) rectangle (2.3838796965159292,4.7481014460444655);
    	\draw[color={black},preaction={fill,color = {green}}] (0.491151,0.952682)--(1.310303,1.148857)--(1.310303,2.08855)--(0.491151,1.892375)--cycle;
	\draw[color={black},preaction={fill,color = {white}}] (1.310303,2.08855)--(0.736727,2.368716)--(-0.082425,2.172542)--(0.491151,1.892375)--cycle;
	\draw[color={black},preaction={fill,color = {gray}}] (0.491151,1.892375)--(-0.082425,2.172542)--(-0.082425,1.232849)--(0.491151,0.952682)--cycle;
	\draw[color={black},preaction={fill,color = {green}}] (0.491151,1.892375)--(1.310303,2.08855)--(1.310303,3.028242)--(0.491151,2.832068)--cycle;
	\draw[color={black},preaction={fill,color = {white}}] (1.310303,3.028242)--(0.736727,3.308409)--(-0.082425,3.112234)--(0.491151,2.832068)--cycle;
	\draw[color={black},preaction={fill,color = {gray}}] (0.491151,2.832068)--(-0.082425,3.112234)--(-0.082425,2.172542)--(0.491151,1.892375)--cycle;
	\draw[color={black},preaction={fill,color = {green}}] (0.491151,2.832068)--(1.310303,3.028242)--(1.310303,3.967935)--(0.491151,3.77176)--cycle;
	\draw[color={black},preaction={fill,color = {white}}] (1.310303,3.967935)--(0.736727,4.248101)--(-0.082425,4.051927)--(0.491151,3.77176)--cycle;
	\draw[color={black},preaction={fill,color = {gray}}] (0.491151,3.77176)--(-0.082425,4.051927)--(-0.082425,3.112234)--(0.491151,2.832068)--cycle;
	\draw[color={black},preaction={fill,color = {green}}] (1.064728,0.672516)--(1.88388,0.868691)--(1.88388,1.808383)--(1.064728,1.612209)--cycle;
	\draw[color={black},preaction={fill,color = {white}}] (1.88388,1.808383)--(1.310303,2.08855)--(0.491151,1.892375)--(1.064728,1.612209)--cycle;
	\draw[color={black},preaction={fill,color = {gray}}] (1.064728,1.612209)--(0.491151,1.892375)--(0.491151,0.952682)--(1.064728,0.672516)--cycle;
	\draw[color={black},preaction={fill,color = {green}}] (1.064728,1.612209)--(1.88388,1.808383)--(1.88388,2.748076)--(1.064728,2.551901)--cycle;
	\draw[color={black},preaction={fill,color = {white}}] (1.88388,2.748076)--(1.310303,3.028242)--(0.491151,2.832068)--(1.064728,2.551901)--cycle;
	\draw[color={black},preaction={fill,color = {gray}}] (1.064728,2.551901)--(0.491151,2.832068)--(0.491151,1.892375)--(1.064728,1.612209)--cycle;
	\draw[color={black},preaction={fill,color = {green}}] (-0.328001,0.756508)--(0.491151,0.952682)--(0.491151,1.892375)--(-0.328001,1.6962)--cycle;
	\draw[color={black},preaction={fill,color = {white}}] (0.491151,1.892375)--(-0.082425,2.172542)--(-0.901577,1.976367)--(-0.328001,1.6962)--cycle;
	\draw[color={black},preaction={fill,color = {gray}}] (-0.328001,1.6962)--(-0.901577,1.976367)--(-0.901577,1.036674)--(-0.328001,0.756508)--cycle;
	\draw[color={black},preaction={fill,color = {green}}] (-0.328001,1.6962)--(0.491151,1.892375)--(0.491151,2.832068)--(-0.328001,2.635893)--cycle;
	\draw[color={black},preaction={fill,color = {white}}] (0.491151,2.832068)--(-0.082425,3.112234)--(-0.901577,2.916059)--(-0.328001,2.635893)--cycle;
	\draw[color={black},preaction={fill,color = {gray}}] (-0.328001,2.635893)--(-0.901577,2.916059)--(-0.901577,1.976367)--(-0.328001,1.6962)--cycle;
	\draw[color={black},preaction={fill,color = {green}}] (-0.328001,2.635893)--(0.491151,2.832068)--(0.491151,3.77176)--(-0.328001,3.575586)--cycle;
	\draw[color={black},preaction={fill,color = {white}}] (0.491151,3.77176)--(-0.082425,4.051927)--(-0.901577,3.855752)--(-0.328001,3.575586)--cycle;
	\draw[color={black},preaction={fill,color = {gray}}] (-0.328001,3.575586)--(-0.901577,3.855752)--(-0.901577,2.916059)--(-0.328001,2.635893)--cycle;
	\draw[color={black},preaction={fill,color = {green}}] (0.245576,0.476341)--(1.064728,0.672516)--(1.064728,1.612209)--(0.245576,1.416034)--cycle;
	\draw[color={black},preaction={fill,color = {white}}] (1.064728,1.612209)--(0.491151,1.892375)--(-0.328001,1.6962)--(0.245576,1.416034)--cycle;
	\draw[color={black},preaction={fill,color = {gray}}] (0.245576,1.416034)--(-0.328001,1.6962)--(-0.328001,0.756508)--(0.245576,0.476341)--cycle;
	\draw[color={black},preaction={fill,color = {green}}] (0.245576,1.416034)--(1.064728,1.612209)--(1.064728,2.551901)--(0.245576,2.355726)--cycle;
	\draw[color={black},preaction={fill,color = {white}}] (1.064728,2.551901)--(0.491151,2.832068)--(-0.328001,2.635893)--(0.245576,2.355726)--cycle;
	\draw[color={black},preaction={fill,color = {gray}}] (0.245576,2.355726)--(-0.328001,2.635893)--(-0.328001,1.6962)--(0.245576,1.416034)--cycle;
	\draw[color={black},preaction={fill,color = {green}}] (-1.147153,0.560333)--(-0.328001,0.756508)--(-0.328001,1.6962)--(-1.147153,1.500026)--cycle;
	\draw[color={black},preaction={fill,color = {white}}] (-0.328001,1.6962)--(-0.901577,1.976367)--(-1.720729,1.780192)--(-1.147153,1.500026)--cycle;
	\draw[color={black},preaction={fill,color = {gray}}] (-1.147153,1.500026)--(-1.720729,1.780192)--(-1.720729,0.840499)--(-1.147153,0.560333)--cycle;
	\draw[color={black},preaction={fill,color = {green}}] (-1.147153,1.500026)--(-0.328001,1.6962)--(-0.328001,2.635893)--(-1.147153,2.439718)--cycle;
	\draw[color={black},preaction={fill,color = {white}}] (-0.328001,2.635893)--(-0.901577,2.916059)--(-1.720729,2.719885)--(-1.147153,2.439718)--cycle;
	\draw[color={black},preaction={fill,color = {gray}}] (-1.147153,2.439718)--(-1.720729,2.719885)--(-1.720729,1.780192)--(-1.147153,1.500026)--cycle;
	\draw[color={black},preaction={fill,color = {green}}] (-1.147153,2.439718)--(-0.328001,2.635893)--(-0.328001,3.575586)--(-1.147153,3.379411)--cycle;
	\draw[color={black},preaction={fill,color = {white}}] (-0.328001,3.575586)--(-0.901577,3.855752)--(-1.720729,3.659577)--(-1.147153,3.379411)--cycle;
	\draw[color={black},preaction={fill,color = {gray}}] (-1.147153,3.379411)--(-1.720729,3.659577)--(-1.720729,2.719885)--(-1.147153,2.439718)--cycle;
	\draw[color={black},preaction={fill,color = {green}}] (-0.573576,0.280166)--(0.245576,0.476341)--(0.245576,1.416034)--(-0.573576,1.219859)--cycle;
	\draw[color={black},preaction={fill,color = {white}}] (0.245576,1.416034)--(-0.328001,1.6962)--(-1.147153,1.500026)--(-0.573576,1.219859)--cycle;
	\draw[color={black},preaction={fill,color = {gray}}] (-0.573576,1.219859)--(-1.147153,1.500026)--(-1.147153,0.560333)--(-0.573576,0.280166)--cycle;
	\draw[color={black},preaction={fill,color = {green}}] (-0.573576,1.219859)--(0.245576,1.416034)--(0.245576,2.355726)--(-0.573576,2.159552)--cycle;
	\draw[color={black},preaction={fill,color = {white}}] (0.245576,2.355726)--(-0.328001,2.635893)--(-1.147153,2.439718)--(-0.573576,2.159552)--cycle;
	\draw[color={black},preaction={fill,color = {gray}}] (-0.573576,2.159552)--(-1.147153,2.439718)--(-1.147153,1.500026)--(-0.573576,1.219859)--cycle;
	\draw[color={black},preaction={fill,color = {green}}] (0,0)--(0.819152,0.196175)--(0.819152,1.135867)--(0,0.939693)--cycle;
	\draw[color={black},preaction={fill,color = {white}}] (0.819152,1.135867)--(0.245576,1.416034)--(-0.573576,1.219859)--(0,0.939693)--cycle;
	\draw[color={black},preaction={fill,color = {gray}}] (0,0.939693)--(-0.573576,1.219859)--(-0.573576,0.280166)--(0,0)--cycle;

\end{tikzpicture}\begin{tikzpicture}[baseline,scale = 0.6]

    \tikzset{
      point/.style={
        thick,
        draw,
        cross out,
        inner sep=0pt,
        minimum width=5pt,
        minimum height=5pt,
      },
    }
    \clip (-3.647152872702092,-0.5) rectangle (3.776608177155967,3.7244170206343323);
    	\draw[color={black},preaction={fill,color = {green}}] (1.88388,0.868691)--(2.703032,1.064865)--(2.703032,2.004558)--(1.88388,1.808383)--cycle;
	\draw[color={black},preaction={fill,color = {white}}] (2.703032,2.004558)--(2.129455,2.284724)--(1.310303,2.08855)--(1.88388,1.808383)--cycle;
	\draw[color={black},preaction={fill,color = {gray}}] (1.88388,1.808383)--(1.310303,2.08855)--(1.310303,1.148857)--(1.88388,0.868691)--cycle;
	\draw[color={black},preaction={fill,color = {green}}] (1.88388,1.808383)--(2.703032,2.004558)--(2.703032,2.944251)--(1.88388,2.748076)--cycle;
	\draw[color={black},preaction={fill,color = {white}}] (2.703032,2.944251)--(2.129455,3.224417)--(1.310303,3.028242)--(1.88388,2.748076)--cycle;
	\draw[color={black},preaction={fill,color = {gray}}] (1.88388,2.748076)--(1.310303,3.028242)--(1.310303,2.08855)--(1.88388,1.808383)--cycle;
	\draw[color={black},preaction={fill,color = {green}}] (2.457456,0.588524)--(3.276608,0.784699)--(3.276608,1.724391)--(2.457456,1.528217)--cycle;
	\draw[color={black},preaction={fill,color = {white}}] (3.276608,1.724391)--(2.703032,2.004558)--(1.88388,1.808383)--(2.457456,1.528217)--cycle;
	\draw[color={black},preaction={fill,color = {gray}}] (2.457456,1.528217)--(1.88388,1.808383)--(1.88388,0.868691)--(2.457456,0.588524)--cycle;
	\draw[color={black},preaction={fill,color = {green}}] (2.457456,1.528217)--(3.276608,1.724391)--(3.276608,2.664084)--(2.457456,2.467909)--cycle;
	\draw[color={black},preaction={fill,color = {white}}] (3.276608,2.664084)--(2.703032,2.944251)--(1.88388,2.748076)--(2.457456,2.467909)--cycle;
	\draw[color={black},preaction={fill,color = {gray}}] (2.457456,2.467909)--(1.88388,2.748076)--(1.88388,1.808383)--(2.457456,1.528217)--cycle;
	\draw[color={black},preaction={fill,color = {green}}] (1.064728,0.672516)--(1.88388,0.868691)--(1.88388,1.808383)--(1.064728,1.612209)--cycle;
	\draw[color={black},preaction={fill,color = {white}}] (1.88388,1.808383)--(1.310303,2.08855)--(0.491151,1.892375)--(1.064728,1.612209)--cycle;
	\draw[color={black},preaction={fill,color = {gray}}] (1.064728,1.612209)--(0.491151,1.892375)--(0.491151,0.952682)--(1.064728,0.672516)--cycle;
	\draw[color={black},preaction={fill,color = {green}}] (1.064728,1.612209)--(1.88388,1.808383)--(1.88388,2.748076)--(1.064728,2.551901)--cycle;
	\draw[color={black},preaction={fill,color = {white}}] (1.88388,2.748076)--(1.310303,3.028242)--(0.491151,2.832068)--(1.064728,2.551901)--cycle;
	\draw[color={black},preaction={fill,color = {gray}}] (1.064728,2.551901)--(0.491151,2.832068)--(0.491151,1.892375)--(1.064728,1.612209)--cycle;
	\draw[color={black},preaction={fill,color = {green}}] (1.638304,0.392349)--(2.457456,0.588524)--(2.457456,1.528217)--(1.638304,1.332042)--cycle;
	\draw[color={black},preaction={fill,color = {white}}] (2.457456,1.528217)--(1.88388,1.808383)--(1.064728,1.612209)--(1.638304,1.332042)--cycle;
	\draw[color={black},preaction={fill,color = {gray}}] (1.638304,1.332042)--(1.064728,1.612209)--(1.064728,0.672516)--(1.638304,0.392349)--cycle;
	\draw[color={black},preaction={fill,color = {green}}] (1.638304,1.332042)--(2.457456,1.528217)--(2.457456,2.467909)--(1.638304,2.271735)--cycle;
	\draw[color={black},preaction={fill,color = {white}}] (2.457456,2.467909)--(1.88388,2.748076)--(1.064728,2.551901)--(1.638304,2.271735)--cycle;
	\draw[color={black},preaction={fill,color = {gray}}] (1.638304,2.271735)--(1.064728,2.551901)--(1.064728,1.612209)--(1.638304,1.332042)--cycle;
	\draw[color={black},preaction={fill,color = {green}}] (0.245576,0.476341)--(1.064728,0.672516)--(1.064728,1.612209)--(0.245576,1.416034)--cycle;
	\draw[color={black},preaction={fill,color = {white}}] (1.064728,1.612209)--(0.491151,1.892375)--(-0.328001,1.6962)--(0.245576,1.416034)--cycle;
	\draw[color={black},preaction={fill,color = {gray}}] (0.245576,1.416034)--(-0.328001,1.6962)--(-0.328001,0.756508)--(0.245576,0.476341)--cycle;
	\draw[color={black},preaction={fill,color = {green}}] (0.245576,1.416034)--(1.064728,1.612209)--(1.064728,2.551901)--(0.245576,2.355726)--cycle;
	\draw[color={black},preaction={fill,color = {white}}] (1.064728,2.551901)--(0.491151,2.832068)--(-0.328001,2.635893)--(0.245576,2.355726)--cycle;
	\draw[color={black},preaction={fill,color = {gray}}] (0.245576,2.355726)--(-0.328001,2.635893)--(-0.328001,1.6962)--(0.245576,1.416034)--cycle;
	\draw[color={black},preaction={fill,color = {green}}] (0.819152,0.196175)--(1.638304,0.392349)--(1.638304,1.332042)--(0.819152,1.135867)--cycle;
	\draw[color={black},preaction={fill,color = {white}}] (1.638304,1.332042)--(1.064728,1.612209)--(0.245576,1.416034)--(0.819152,1.135867)--cycle;
	\draw[color={black},preaction={fill,color = {gray}}] (0.819152,1.135867)--(0.245576,1.416034)--(0.245576,0.476341)--(0.819152,0.196175)--cycle;
	\draw[color={black},preaction={fill,color = {green}}] (0.819152,1.135867)--(1.638304,1.332042)--(1.638304,2.271735)--(0.819152,2.07556)--cycle;
	\draw[color={black},preaction={fill,color = {white}}] (1.638304,2.271735)--(1.064728,2.551901)--(0.245576,2.355726)--(0.819152,2.07556)--cycle;
	\draw[color={black},preaction={fill,color = {gray}}] (0.819152,2.07556)--(0.245576,2.355726)--(0.245576,1.416034)--(0.819152,1.135867)--cycle;
	\draw[color={black},preaction={fill,color = {green}}] (-0.573576,0.280166)--(0.245576,0.476341)--(0.245576,1.416034)--(-0.573576,1.219859)--cycle;
	\draw[color={black},preaction={fill,color = {white}}] (0.245576,1.416034)--(-0.328001,1.6962)--(-1.147153,1.500026)--(-0.573576,1.219859)--cycle;
	\draw[color={black},preaction={fill,color = {gray}}] (-0.573576,1.219859)--(-1.147153,1.500026)--(-1.147153,0.560333)--(-0.573576,0.280166)--cycle;
	\draw[color={black},preaction={fill,color = {green}}] (-0.573576,1.219859)--(0.245576,1.416034)--(0.245576,2.355726)--(-0.573576,2.159552)--cycle;
	\draw[color={black},preaction={fill,color = {white}}] (0.245576,2.355726)--(-0.328001,2.635893)--(-1.147153,2.439718)--(-0.573576,2.159552)--cycle;
	\draw[color={black},preaction={fill,color = {gray}}] (-0.573576,2.159552)--(-1.147153,2.439718)--(-1.147153,1.500026)--(-0.573576,1.219859)--cycle;
	\draw[color={black},preaction={fill,color = {green}}] (0,0)--(0.819152,0.196175)--(0.819152,1.135867)--(0,0.939693)--cycle;
	\draw[color={black},preaction={fill,color = {white}}] (0.819152,1.135867)--(0.245576,1.416034)--(-0.573576,1.219859)--(0,0.939693)--cycle;
	\draw[color={black},preaction={fill,color = {gray}}] (0,0.939693)--(-0.573576,1.219859)--(-0.573576,0.280166)--(0,0)--cycle;
	\draw[color={black},preaction={fill,color = {green}}] (0,0.939693)--(0.819152,1.135867)--(0.819152,2.07556)--(0,1.879385)--cycle;
	\draw[color={black},preaction={fill,color = {white}}] (0.819152,2.07556)--(0.245576,2.355726)--(-0.573576,2.159552)--(0,1.879385)--cycle;
	\draw[color={black},preaction={fill,color = {gray}}] (0,1.879385)--(-0.573576,2.159552)--(-0.573576,1.219859)--(0,0.939693)--cycle;

\end{tikzpicture}\\Quel est celui qui a le plus petit volume (en nombre de cubes identiques) ?
\end{EXO}

% @see : https://coopmaths.fr/alea?uuid=e332d&id=6M3C-flash1&n=1&d=10&s=false&alea=X9Tn&cd=1&cols=1
\begin{EXO}{}{6M3C-flash1}

 \begin{tikzpicture}[baseline]

    \tikzset{
      point/.style={
        thick,
        draw,
        cross out,
        inner sep=0pt,
        minimum width=5pt,
        minimum height=5pt,
      },
    }
    \clip (-0.5,-0.5) rectangle (4.8,3.5);
    	\draw[color={black},preaction={fill,color = {lightgray}, opacity = 0.5}] (0,0)--(0.75,0)--(0.75,0.75)--(0,0.75)--cycle;
	\draw[color={black},preaction={fill,color = {lightgray}, opacity = 0.5}] (0,0.75)--(0.75,0.75)--(0.75,1.5)--(0,1.5)--cycle;
	\draw[color={black},preaction={fill,color = {lightgray}, opacity = 0.5}] (0,1.5)--(0.75,1.5)--(0.75,2.25)--(0,2.25)--cycle;
	\draw[color={black},preaction={fill,color = {white}, opacity = 0.5}] (0,2.25)--(0.75,2.25)--(1.07,2.44)--(0.32,2.44)--cycle;
	\draw[color={black},preaction={fill,color = {white}, opacity = 0.5}] (0.32,2.44)--(1.07,2.44)--(1.4,2.63)--(0.65,2.63)--cycle;
	\draw[color={black},preaction={fill,color = {white}, opacity = 0.5}] (0.65,2.63)--(1.4,2.63)--(1.72,2.81)--(0.97,2.81)--cycle;
	\draw[color={black},preaction={fill,color = {white}, opacity = 0.5}] (0.97,2.81)--(1.72,2.81)--(2.05,3)--(1.3,3)--cycle;
	\draw[color={black},preaction={fill,color = {lightgray}, opacity = 0.5}] (0.75,0)--(1.5,0)--(1.5,0.75)--(0.75,0.75)--cycle;
	\draw[color={black},preaction={fill,color = {lightgray}, opacity = 0.5}] (0.75,0.75)--(1.5,0.75)--(1.5,1.5)--(0.75,1.5)--cycle;
	\draw[color={black},preaction={fill,color = {lightgray}, opacity = 0.5}] (0.75,1.5)--(1.5,1.5)--(1.5,2.25)--(0.75,2.25)--cycle;
	\draw[color={black},preaction={fill,color = {white}, opacity = 0.5}] (0.75,2.25)--(1.5,2.25)--(1.82,2.44)--(1.07,2.44)--cycle;
	\draw[color={black},preaction={fill,color = {white}, opacity = 0.5}] (1.07,2.44)--(1.82,2.44)--(2.15,2.63)--(1.4,2.63)--cycle;
	\draw[color={black},preaction={fill,color = {white}, opacity = 0.5}] (1.4,2.63)--(2.15,2.63)--(2.47,2.81)--(1.72,2.81)--cycle;
	\draw[color={black},preaction={fill,color = {white}, opacity = 0.5}] (1.72,2.81)--(2.47,2.81)--(2.8,3)--(2.05,3)--cycle;
	\draw[color={black},preaction={fill,color = {lightgray}, opacity = 0.5}] (1.5,0)--(2.25,0)--(2.25,0.75)--(1.5,0.75)--cycle;
	\draw[color={black},preaction={fill,color = {lightgray}, opacity = 0.5}] (1.5,0.75)--(2.25,0.75)--(2.25,1.5)--(1.5,1.5)--cycle;
	\draw[color={black},preaction={fill,color = {lightgray}, opacity = 0.5}] (1.5,1.5)--(2.25,1.5)--(2.25,2.25)--(1.5,2.25)--cycle;
	\draw[color={black},preaction={fill,color = {white}, opacity = 0.5}] (1.5,2.25)--(2.25,2.25)--(2.57,2.44)--(1.82,2.44)--cycle;
	\draw[color={black},preaction={fill,color = {white}, opacity = 0.5}] (1.82,2.44)--(2.57,2.44)--(2.9,2.63)--(2.15,2.63)--cycle;
	\draw[color={black},preaction={fill,color = {white}, opacity = 0.5}] (2.15,2.63)--(2.9,2.63)--(3.22,2.81)--(2.47,2.81)--cycle;
	\draw[color={black},preaction={fill,color = {white}, opacity = 0.5}] (2.47,2.81)--(3.22,2.81)--(3.55,3)--(2.8,3)--cycle;
	\draw[color={black},preaction={fill,color = {lightgray}, opacity = 0.5}] (2.25,0)--(3,0)--(3,0.75)--(2.25,0.75)--cycle;
	\draw[color={black},preaction={fill,color = {darkgray}, opacity = 0.5}] (3,0)--(3.32,0.19)--(3.32,0.94)--(3,0.75)--cycle;
	\draw[color={black},preaction={fill,color = {lightgray}, opacity = 0.5}] (2.25,0.75)--(3,0.75)--(3,1.5)--(2.25,1.5)--cycle;
	\draw[color={black},preaction={fill,color = {darkgray}, opacity = 0.5}] (3,0.75)--(3.32,0.94)--(3.32,1.69)--(3,1.5)--cycle;
	\draw[color={black},preaction={fill,color = {lightgray}, opacity = 0.5}] (2.25,1.5)--(3,1.5)--(3,2.25)--(2.25,2.25)--cycle;
	\draw[color={black},preaction={fill,color = {darkgray}, opacity = 0.5}] (3,1.5)--(3.32,1.69)--(3.32,2.44)--(3,2.25)--cycle;
	\draw[color={black},preaction={fill,color = {white}, opacity = 0.5}] (2.25,2.25)--(3,2.25)--(3.32,2.44)--(2.57,2.44)--cycle;
	\draw[color={black},preaction={fill,color = {darkgray}, opacity = 0.5}] (3.32,0.19)--(3.65,0.37)--(3.65,1.13)--(3.32,0.94)--cycle;
	\draw[color={black},preaction={fill,color = {darkgray}, opacity = 0.5}] (3.32,0.94)--(3.65,1.13)--(3.65,1.88)--(3.32,1.69)--cycle;
	\draw[color={black},preaction={fill,color = {darkgray}, opacity = 0.5}] (3.32,1.69)--(3.65,1.88)--(3.65,2.63)--(3.32,2.44)--cycle;
	\draw[color={black},preaction={fill,color = {white}, opacity = 0.5}] (2.57,2.44)--(3.32,2.44)--(3.65,2.63)--(2.9,2.63)--cycle;
	\draw[color={black},preaction={fill,color = {darkgray}, opacity = 0.5}] (3.65,0.37)--(3.97,0.56)--(3.97,1.31)--(3.65,1.13)--cycle;
	\draw[color={black},preaction={fill,color = {darkgray}, opacity = 0.5}] (3.65,1.13)--(3.97,1.31)--(3.97,2.06)--(3.65,1.88)--cycle;
	\draw[color={black},preaction={fill,color = {darkgray}, opacity = 0.5}] (3.65,1.88)--(3.97,2.06)--(3.97,2.81)--(3.65,2.63)--cycle;
	\draw[color={black},preaction={fill,color = {white}, opacity = 0.5}] (2.9,2.63)--(3.65,2.63)--(3.97,2.81)--(3.22,2.81)--cycle;
	\draw[color={black},preaction={fill,color = {darkgray}, opacity = 0.5}] (3.97,0.56)--(4.3,0.75)--(4.3,1.5)--(3.97,1.31)--cycle;
	\draw[color={black},preaction={fill,color = {darkgray}, opacity = 0.5}] (3.97,1.31)--(4.3,1.5)--(4.3,2.25)--(3.97,2.06)--cycle;
	\draw[color={black},preaction={fill,color = {darkgray}, opacity = 0.5}] (3.97,2.06)--(4.3,2.25)--(4.3,3)--(3.97,2.81)--cycle;
	\draw[color={black},preaction={fill,color = {white}, opacity = 0.5}] (3.22,2.81)--(3.97,2.81)--(4.3,3)--(3.55,3)--cycle;

\end{tikzpicture}\\

\end{EXO}


\clearpage

\begin{Correction}
\begin{EXO}{}{}

 Prisme droit avec une base ayant $4$ sommets.
\end{EXO}

\begin{EXO}{}{}

 Pyramide avec une base ayant $5$ sommets.
\end{EXO}

\begin{EXO}{}{}

 Cône de révolution
\end{EXO}

\begin{EXO}{}{}

 Pavé droit.
\end{EXO}

\begin{EXO}{}{}

 Cube.
\end{EXO}

\begin{EXO}{}{}

 Sphère.
\end{EXO}

\begin{EXO}{}{}

 \begin{tikzpicture}[baseline,scale = 0.6]

    \tikzset{
      point/.style={
        thick,
        draw,
        cross out,
        inner sep=0pt,
        minimum width=5pt,
        minimum height=5pt,
      },
    }
    \clip (-2.220729309053138,-0.5) rectangle (5.660487873671896,5.532800225920509);
    	\draw[color={black},preaction={fill,color = {green}}] (3.767759,1.737381)--(4.586911,1.933556)--(4.586911,2.873248)--(3.767759,2.677074)--cycle;
	\draw[color={black},preaction={fill,color = {white}}] (4.586911,2.873248)--(4.013335,3.153415)--(3.194183,2.95724)--(3.767759,2.677074)--cycle;
	\draw[color={black},preaction={fill,color = {gray}}] (3.767759,2.677074)--(3.194183,2.95724)--(3.194183,2.017548)--(3.767759,1.737381)--cycle;
	\draw[color={black},preaction={fill,color = {green}}] (3.767759,2.677074)--(4.586911,2.873248)--(4.586911,3.812941)--(3.767759,3.616766)--cycle;
	\draw[color={black},preaction={fill,color = {white}}] (4.586911,3.812941)--(4.013335,4.093108)--(3.194183,3.896933)--(3.767759,3.616766)--cycle;
	\draw[color={black},preaction={fill,color = {gray}}] (3.767759,3.616766)--(3.194183,3.896933)--(3.194183,2.95724)--(3.767759,2.677074)--cycle;
	\draw[color={black},preaction={fill,color = {green}}] (3.767759,3.616766)--(4.586911,3.812941)--(4.586911,4.752634)--(3.767759,4.556459)--cycle;
	\draw[color={black},preaction={fill,color = {white}}] (4.586911,4.752634)--(4.013335,5.0328)--(3.194183,4.836626)--(3.767759,4.556459)--cycle;
	\draw[color={black},preaction={fill,color = {gray}}] (3.767759,4.556459)--(3.194183,4.836626)--(3.194183,3.896933)--(3.767759,3.616766)--cycle;
	\draw[color={black},preaction={fill,color = {green}}] (4.341336,1.457215)--(5.160488,1.653389)--(5.160488,2.593082)--(4.341336,2.396907)--cycle;
	\draw[color={black},preaction={fill,color = {white}}] (5.160488,2.593082)--(4.586911,2.873248)--(3.767759,2.677074)--(4.341336,2.396907)--cycle;
	\draw[color={black},preaction={fill,color = {gray}}] (4.341336,2.396907)--(3.767759,2.677074)--(3.767759,1.737381)--(4.341336,1.457215)--cycle;
	\draw[color={black},preaction={fill,color = {green}}] (4.341336,2.396907)--(5.160488,2.593082)--(5.160488,3.532775)--(4.341336,3.3366)--cycle;
	\draw[color={black},preaction={fill,color = {white}}] (5.160488,3.532775)--(4.586911,3.812941)--(3.767759,3.616766)--(4.341336,3.3366)--cycle;
	\draw[color={black},preaction={fill,color = {gray}}] (4.341336,3.3366)--(3.767759,3.616766)--(3.767759,2.677074)--(4.341336,2.396907)--cycle;
	\draw[color={black},preaction={fill,color = {green}}] (1.310303,1.148857)--(2.129455,1.345032)--(2.129455,2.284724)--(1.310303,2.08855)--cycle;
	\draw[color={black},preaction={fill,color = {white}}] (2.129455,2.284724)--(1.555879,2.564891)--(0.736727,2.368716)--(1.310303,2.08855)--cycle;
	\draw[color={black},preaction={fill,color = {gray}}] (1.310303,2.08855)--(0.736727,2.368716)--(0.736727,1.429024)--(1.310303,1.148857)--cycle;
	\draw[color={black},preaction={fill,color = {green}}] (1.310303,2.08855)--(2.129455,2.284724)--(2.129455,3.224417)--(1.310303,3.028242)--cycle;
	\draw[color={black},preaction={fill,color = {white}}] (2.129455,3.224417)--(1.555879,3.504584)--(0.736727,3.308409)--(1.310303,3.028242)--cycle;
	\draw[color={black},preaction={fill,color = {gray}}] (1.310303,3.028242)--(0.736727,3.308409)--(0.736727,2.368716)--(1.310303,2.08855)--cycle;
	\draw[color={black},preaction={fill,color = {green}}] (1.310303,3.028242)--(2.129455,3.224417)--(2.129455,4.16411)--(1.310303,3.967935)--cycle;
	\draw[color={black},preaction={fill,color = {white}}] (2.129455,4.16411)--(1.555879,4.444276)--(0.736727,4.248101)--(1.310303,3.967935)--cycle;
	\draw[color={black},preaction={fill,color = {gray}}] (1.310303,3.967935)--(0.736727,4.248101)--(0.736727,3.308409)--(1.310303,3.028242)--cycle;
	\draw[color={black},preaction={fill,color = {green}}] (1.88388,0.868691)--(2.703032,1.064865)--(2.703032,2.004558)--(1.88388,1.808383)--cycle;
	\draw[color={black},preaction={fill,color = {white}}] (2.703032,2.004558)--(2.129455,2.284724)--(1.310303,2.08855)--(1.88388,1.808383)--cycle;
	\draw[color={black},preaction={fill,color = {gray}}] (1.88388,1.808383)--(1.310303,2.08855)--(1.310303,1.148857)--(1.88388,0.868691)--cycle;
	\draw[color={black},preaction={fill,color = {green}}] (1.88388,1.808383)--(2.703032,2.004558)--(2.703032,2.944251)--(1.88388,2.748076)--cycle;
	\draw[color={black},preaction={fill,color = {white}}] (2.703032,2.944251)--(2.129455,3.224417)--(1.310303,3.028242)--(1.88388,2.748076)--cycle;
	\draw[color={black},preaction={fill,color = {gray}}] (1.88388,2.748076)--(1.310303,3.028242)--(1.310303,2.08855)--(1.88388,1.808383)--cycle;
	\draw[color={black},preaction={fill,color = {green}}] (-1.147153,0.560333)--(-0.328001,0.756508)--(-0.328001,1.6962)--(-1.147153,1.500026)--cycle;
	\draw[color={black},preaction={fill,color = {white}}] (-0.328001,1.6962)--(-0.901577,1.976367)--(-1.720729,1.780192)--(-1.147153,1.500026)--cycle;
	\draw[color={black},preaction={fill,color = {gray}}] (-1.147153,1.500026)--(-1.720729,1.780192)--(-1.720729,0.840499)--(-1.147153,0.560333)--cycle;
	\draw[color={black},preaction={fill,color = {green}}] (-1.147153,1.500026)--(-0.328001,1.6962)--(-0.328001,2.635893)--(-1.147153,2.439718)--cycle;
	\draw[color={black},preaction={fill,color = {white}}] (-0.328001,2.635893)--(-0.901577,2.916059)--(-1.720729,2.719885)--(-1.147153,2.439718)--cycle;
	\draw[color={black},preaction={fill,color = {gray}}] (-1.147153,2.439718)--(-1.720729,2.719885)--(-1.720729,1.780192)--(-1.147153,1.500026)--cycle;
	\draw[color={black},preaction={fill,color = {green}}] (-1.147153,2.439718)--(-0.328001,2.635893)--(-0.328001,3.575586)--(-1.147153,3.379411)--cycle;
	\draw[color={black},preaction={fill,color = {white}}] (-0.328001,3.575586)--(-0.901577,3.855752)--(-1.720729,3.659577)--(-1.147153,3.379411)--cycle;
	\draw[color={black},preaction={fill,color = {gray}}] (-1.147153,3.379411)--(-1.720729,3.659577)--(-1.720729,2.719885)--(-1.147153,2.439718)--cycle;
	\draw[color={black},preaction={fill,color = {green}}] (-0.573576,0.280166)--(0.245576,0.476341)--(0.245576,1.416034)--(-0.573576,1.219859)--cycle;
	\draw[color={black},preaction={fill,color = {white}}] (0.245576,1.416034)--(-0.328001,1.6962)--(-1.147153,1.500026)--(-0.573576,1.219859)--cycle;
	\draw[color={black},preaction={fill,color = {gray}}] (-0.573576,1.219859)--(-1.147153,1.500026)--(-1.147153,0.560333)--(-0.573576,0.280166)--cycle;
	\draw[color={black},preaction={fill,color = {green}}] (-0.573576,1.219859)--(0.245576,1.416034)--(0.245576,2.355726)--(-0.573576,2.159552)--cycle;
	\draw[color={black},preaction={fill,color = {white}}] (0.245576,2.355726)--(-0.328001,2.635893)--(-1.147153,2.439718)--(-0.573576,2.159552)--cycle;
	\draw[color={black},preaction={fill,color = {gray}}] (-0.573576,2.159552)--(-1.147153,2.439718)--(-1.147153,1.500026)--(-0.573576,1.219859)--cycle;
	\draw[color={black},preaction={fill,color = {green}}] (0,0)--(0.819152,0.196175)--(0.819152,1.135867)--(0,0.939693)--cycle;
	\draw[color={black},preaction={fill,color = {white}}] (0.819152,1.135867)--(0.245576,1.416034)--(-0.573576,1.219859)--(0,0.939693)--cycle;
	\draw[color={black},preaction={fill,color = {gray}}] (0,0.939693)--(-0.573576,1.219859)--(-0.573576,0.280166)--(0,0)--cycle;

\end{tikzpicture}\begin{tikzpicture}[baseline,scale = 0.6]

    \tikzset{
      point/.style={
        thick,
        draw,
        cross out,
        inner sep=0pt,
        minimum width=5pt,
        minimum height=5pt,
      },
    }
    \clip (-3.647152872702092,-0.5) rectangle (8.691520442889917,4.901465190448398);
    	\draw[color={black},preaction={fill,color = {green}}] (6.798792,2.045739)--(7.617944,2.241913)--(7.617944,3.181606)--(6.798792,2.985431)--cycle;
	\draw[color={black},preaction={fill,color = {white}}] (7.617944,3.181606)--(7.044368,3.461773)--(6.225216,3.265598)--(6.798792,2.985431)--cycle;
	\draw[color={black},preaction={fill,color = {gray}}] (6.798792,2.985431)--(6.225216,3.265598)--(6.225216,2.325905)--(6.798792,2.045739)--cycle;
	\draw[color={black},preaction={fill,color = {green}}] (6.798792,2.985431)--(7.617944,3.181606)--(7.617944,4.121299)--(6.798792,3.925124)--cycle;
	\draw[color={black},preaction={fill,color = {white}}] (7.617944,4.121299)--(7.044368,4.401465)--(6.225216,4.20529)--(6.798792,3.925124)--cycle;
	\draw[color={black},preaction={fill,color = {gray}}] (6.798792,3.925124)--(6.225216,4.20529)--(6.225216,3.265598)--(6.798792,2.985431)--cycle;
	\draw[color={black},preaction={fill,color = {green}}] (7.372368,1.765572)--(8.19152,1.961747)--(8.19152,2.90144)--(7.372368,2.705265)--cycle;
	\draw[color={black},preaction={fill,color = {white}}] (8.19152,2.90144)--(7.617944,3.181606)--(6.798792,2.985431)--(7.372368,2.705265)--cycle;
	\draw[color={black},preaction={fill,color = {gray}}] (7.372368,2.705265)--(6.798792,2.985431)--(6.798792,2.045739)--(7.372368,1.765572)--cycle;
	\draw[color={black},preaction={fill,color = {green}}] (7.372368,2.705265)--(8.19152,2.90144)--(8.19152,3.841132)--(7.372368,3.644957)--cycle;
	\draw[color={black},preaction={fill,color = {white}}] (8.19152,3.841132)--(7.617944,4.121299)--(6.798792,3.925124)--(7.372368,3.644957)--cycle;
	\draw[color={black},preaction={fill,color = {gray}}] (7.372368,3.644957)--(6.798792,3.925124)--(6.798792,2.985431)--(7.372368,2.705265)--cycle;
	\draw[color={black},preaction={fill,color = {green}}] (4.341336,1.457215)--(5.160488,1.653389)--(5.160488,2.593082)--(4.341336,2.396907)--cycle;
	\draw[color={black},preaction={fill,color = {white}}] (5.160488,2.593082)--(4.586911,2.873248)--(3.767759,2.677074)--(4.341336,2.396907)--cycle;
	\draw[color={black},preaction={fill,color = {gray}}] (4.341336,2.396907)--(3.767759,2.677074)--(3.767759,1.737381)--(4.341336,1.457215)--cycle;
	\draw[color={black},preaction={fill,color = {green}}] (4.341336,2.396907)--(5.160488,2.593082)--(5.160488,3.532775)--(4.341336,3.3366)--cycle;
	\draw[color={black},preaction={fill,color = {white}}] (5.160488,3.532775)--(4.586911,3.812941)--(3.767759,3.616766)--(4.341336,3.3366)--cycle;
	\draw[color={black},preaction={fill,color = {gray}}] (4.341336,3.3366)--(3.767759,3.616766)--(3.767759,2.677074)--(4.341336,2.396907)--cycle;
	\draw[color={black},preaction={fill,color = {green}}] (4.914912,1.177048)--(5.734064,1.373223)--(5.734064,2.312915)--(4.914912,2.116741)--cycle;
	\draw[color={black},preaction={fill,color = {white}}] (5.734064,2.312915)--(5.160488,2.593082)--(4.341336,2.396907)--(4.914912,2.116741)--cycle;
	\draw[color={black},preaction={fill,color = {gray}}] (4.914912,2.116741)--(4.341336,2.396907)--(4.341336,1.457215)--(4.914912,1.177048)--cycle;
	\draw[color={black},preaction={fill,color = {green}}] (4.914912,2.116741)--(5.734064,2.312915)--(5.734064,3.252608)--(4.914912,3.056433)--cycle;
	\draw[color={black},preaction={fill,color = {white}}] (5.734064,3.252608)--(5.160488,3.532775)--(4.341336,3.3366)--(4.914912,3.056433)--cycle;
	\draw[color={black},preaction={fill,color = {gray}}] (4.914912,3.056433)--(4.341336,3.3366)--(4.341336,2.396907)--(4.914912,2.116741)--cycle;
	\draw[color={black},preaction={fill,color = {green}}] (1.88388,0.868691)--(2.703032,1.064865)--(2.703032,2.004558)--(1.88388,1.808383)--cycle;
	\draw[color={black},preaction={fill,color = {white}}] (2.703032,2.004558)--(2.129455,2.284724)--(1.310303,2.08855)--(1.88388,1.808383)--cycle;
	\draw[color={black},preaction={fill,color = {gray}}] (1.88388,1.808383)--(1.310303,2.08855)--(1.310303,1.148857)--(1.88388,0.868691)--cycle;
	\draw[color={black},preaction={fill,color = {green}}] (1.88388,1.808383)--(2.703032,2.004558)--(2.703032,2.944251)--(1.88388,2.748076)--cycle;
	\draw[color={black},preaction={fill,color = {white}}] (2.703032,2.944251)--(2.129455,3.224417)--(1.310303,3.028242)--(1.88388,2.748076)--cycle;
	\draw[color={black},preaction={fill,color = {gray}}] (1.88388,2.748076)--(1.310303,3.028242)--(1.310303,2.08855)--(1.88388,1.808383)--cycle;
	\draw[color={black},preaction={fill,color = {green}}] (2.457456,0.588524)--(3.276608,0.784699)--(3.276608,1.724391)--(2.457456,1.528217)--cycle;
	\draw[color={black},preaction={fill,color = {white}}] (3.276608,1.724391)--(2.703032,2.004558)--(1.88388,1.808383)--(2.457456,1.528217)--cycle;
	\draw[color={black},preaction={fill,color = {gray}}] (2.457456,1.528217)--(1.88388,1.808383)--(1.88388,0.868691)--(2.457456,0.588524)--cycle;
	\draw[color={black},preaction={fill,color = {green}}] (2.457456,1.528217)--(3.276608,1.724391)--(3.276608,2.664084)--(2.457456,2.467909)--cycle;
	\draw[color={black},preaction={fill,color = {white}}] (3.276608,2.664084)--(2.703032,2.944251)--(1.88388,2.748076)--(2.457456,2.467909)--cycle;
	\draw[color={black},preaction={fill,color = {gray}}] (2.457456,2.467909)--(1.88388,2.748076)--(1.88388,1.808383)--(2.457456,1.528217)--cycle;
	\draw[color={black},preaction={fill,color = {green}}] (-0.573576,0.280166)--(0.245576,0.476341)--(0.245576,1.416034)--(-0.573576,1.219859)--cycle;
	\draw[color={black},preaction={fill,color = {white}}] (0.245576,1.416034)--(-0.328001,1.6962)--(-1.147153,1.500026)--(-0.573576,1.219859)--cycle;
	\draw[color={black},preaction={fill,color = {gray}}] (-0.573576,1.219859)--(-1.147153,1.500026)--(-1.147153,0.560333)--(-0.573576,0.280166)--cycle;
	\draw[color={black},preaction={fill,color = {green}}] (-0.573576,1.219859)--(0.245576,1.416034)--(0.245576,2.355726)--(-0.573576,2.159552)--cycle;
	\draw[color={black},preaction={fill,color = {white}}] (0.245576,2.355726)--(-0.328001,2.635893)--(-1.147153,2.439718)--(-0.573576,2.159552)--cycle;
	\draw[color={black},preaction={fill,color = {gray}}] (-0.573576,2.159552)--(-1.147153,2.439718)--(-1.147153,1.500026)--(-0.573576,1.219859)--cycle;
	\draw[color={black},preaction={fill,color = {green}}] (0,0)--(0.819152,0.196175)--(0.819152,1.135867)--(0,0.939693)--cycle;
	\draw[color={black},preaction={fill,color = {white}}] (0.819152,1.135867)--(0.245576,1.416034)--(-0.573576,1.219859)--(0,0.939693)--cycle;
	\draw[color={black},preaction={fill,color = {gray}}] (0,0.939693)--(-0.573576,1.219859)--(-0.573576,0.280166)--(0,0)--cycle;
	\draw[color={black},preaction={fill,color = {green}}] (0,0.939693)--(0.819152,1.135867)--(0.819152,2.07556)--(0,1.879385)--cycle;
	\draw[color={black},preaction={fill,color = {white}}] (0.819152,2.07556)--(0.245576,2.355726)--(-0.573576,2.159552)--(0,1.879385)--cycle;
	\draw[color={black},preaction={fill,color = {gray}}] (0,1.879385)--(-0.573576,2.159552)--(-0.573576,1.219859)--(0,0.939693)--cycle;

\end{tikzpicture}\\Le premier solide est un empilement de 16 cubes.\\Le deuxième solide est un pavé droit de dimensions $4 \times 2 \times 2$, son volume est de $16$.\\
\end{EXO}

\begin{EXO}{}{}

 Le volume de ce pavé droit est : $4\times 4\times 3={\color[HTML]{F15929}\boldsymbol{48}}$ cubes
\end{EXO}

\clearpage
\end{Correction}
\end{document}

% Local Variables:
% TeX-engine: luatex
% End: